\documentclass[10pt,a4paper]{amsart}
\usepackage{amsmath,amssymb,amsthm,graphicx}
% Retained minimal preamble for the Stacks Project chapter varieties.tex.
% Verbatim excerpt of the chapter preamble: Xy-pic setup (native preamble lines 14-18), the native theorem
% declarations with the mathematical macros (native preamble lines 154-217), and the project own no-op
% definitions of the reference/slogan/history display environments (native wrapper.tex lines 1-3) that the
% counted statement uses for its slogan. Copied mechanically; nothing added, removed or redefined.
% For commutative diagrams we use Xy-pic
\usepackage[all]{xy}

% We use 2cell for 2-commutative diagrams.
\xyoption{2cell}
% Theorem environments.
%
\theoremstyle{plain}
\newtheorem{theorem}[subsection]{Theorem}
\newtheorem{proposition}[subsection]{Proposition}
\newtheorem{lemma}[subsection]{Lemma}

\theoremstyle{definition}
\newtheorem{definition}[subsection]{Definition}
\newtheorem{example}[subsection]{Example}
\newtheorem{exercise}[subsection]{Exercise}
\newtheorem{situation}[subsection]{Situation}

\theoremstyle{remark}
\newtheorem{remark}[subsection]{Remark}
\newtheorem{remarks}[subsection]{Remarks}

\numberwithin{equation}{subsection}

% Macros
%
\def\lim{\mathop{\mathrm{lim}}\nolimits}
\def\colim{\mathop{\mathrm{colim}}\nolimits}
\def\Spec{\mathop{\mathrm{Spec}}}
\def\Hom{\mathop{\mathrm{Hom}}\nolimits}
\def\Ext{\mathop{\mathrm{Ext}}\nolimits}
\def\SheafHom{\mathop{\mathcal{H}\!\mathit{om}}\nolimits}
\def\SheafExt{\mathop{\mathcal{E}\!\mathit{xt}}\nolimits}
\def\Sch{\mathit{Sch}}
\def\Mor{\mathop{\mathrm{Mor}}\nolimits}
\def\Ob{\mathop{\mathrm{Ob}}\nolimits}
\def\Sh{\mathop{\mathit{Sh}}\nolimits}
\def\NL{\mathop{N\!L}\nolimits}
\def\CH{\mathop{\mathrm{CH}}\nolimits}
\def\proetale{{pro\text{-}\acute{e}tale}}
\def\etale{{\acute{e}tale}}
\def\QCoh{\mathit{QCoh}}
\def\Ker{\mathop{\mathrm{Ker}}}
\def\Im{\mathop{\mathrm{Im}}}
\def\Coker{\mathop{\mathrm{Coker}}}
\def\Coim{\mathop{\mathrm{Coim}}}

% Boxtimes
%
\DeclareMathSymbol{\boxtimes}{\mathbin}{AMSa}{"02}

%
% Macros for moduli stacks/spaces
%
\def\QCohstack{\mathcal{QC}\!\mathit{oh}}
\def\Cohstack{\mathcal{C}\!\mathit{oh}}
\def\Spacesstack{\mathcal{S}\!\mathit{paces}}
\def\Quotfunctor{\mathrm{Quot}}
\def\Hilbfunctor{\mathrm{Hilb}}
\def\Curvesstack{\mathcal{C}\!\mathit{urves}}
\def\Polarizedstack{\mathcal{P}\!\mathit{olarized}}
\def\Complexesstack{\mathcal{C}\!\mathit{omplexes}}
% \Pic is the operator that assigns to X its picard group, usage \Pic(X)
% \Picardstack_{X/B} denotes the Picard stack of X over B
% \Picardfunctor_{X/B} denotes the Picard functor of X over B
\def\Pic{\mathop{\mathrm{Pic}}\nolimits}
\def\Picardstack{\mathcal{P}\!\mathit{ic}}
\def\Picardfunctor{\mathrm{Pic}}
\def\Deformationcategory{\mathcal{D}\!\mathit{ef}}
\newenvironment{reference}{}{}
\newenvironment{slogan}{}{}
\newenvironment{history}{}{}
\usepackage{cleveref}
\crefname{originallemma}{original lemma}{original lemmas}
\Crefname{originallemma}{Original lemma}{Original lemmas}
\makeatletter
\newlabel{morphisms-lemma-dimension-fibre-at-a-point}{{02FX}{source}{}{sourcepointer-1}{}}
\newlabel{morphisms-lemma-dimension-fibre-at-a-point@cref}{{[originallemma][0][]02FX}{[1][0][]source}{}{}{}}
\makeatother
\crefname{originallemma}{original lemma}{original lemmas}
\Crefname{originallemma}{Original lemma}{Original lemmas}
\makeatletter
\newlabel{morphisms-lemma-dimension-formula}{{02JU}{source}{}{sourcepointer-2}{}}
\newlabel{morphisms-lemma-dimension-formula@cref}{{[originallemma][0][]02JU}{[1][0][]source}{}{}{}}
\makeatother
\crefname{originallemma}{original lemma}{original lemmas}
\Crefname{originallemma}{Original lemma}{Original lemmas}
\makeatletter
\newlabel{morphisms-lemma-openness-bounded-dimension-fibres}{{02FZ}{source}{}{sourcepointer-3}{}}
\newlabel{morphisms-lemma-openness-bounded-dimension-fibres@cref}{{[originallemma][0][]02FZ}{[1][0][]source}{}{}{}}
\makeatother
\crefname{originallemma}{original lemma}{original lemmas}
\Crefname{originallemma}{Original lemma}{Original lemmas}
\makeatletter
\newlabel{schemes-lemma-specialize-points}{{01J7}{source}{}{sourcepointer-4}{}}
\newlabel{schemes-lemma-specialize-points@cref}{{[originallemma][0][]01J7}{[1][0][]source}{}{}{}}
\makeatother
\title{Dimension of a fibre over a one-dimensional local ring from the dimension at generic points}
\author{The Stacks Project Authors}
\date{stacks-project:varieties.tex:local-snapshot-20261011; GNU FDL 1.2 or later}
\begin{document}
\maketitle
\noindent\textit{Editorial note.} All delivered mathematical text is the author's own, copied from the pinned Stacks Project chapter file varieties.tex; no statement or proof step was rewritten, reordered or removed.\par\medskip
\noindent\textit{Editorial note.} The proof is complete as written and ends with the two matching inequalities; no step is referred to the reader.\par\medskip
\noindent\textit{Editorial note.} The counted spans carry 5 navigation stubs and no bibliography citation.\par\medskip
\noindent\textit{Editorial note.} The preamble is the chapter's retained minimal preamble (sources/stacks-varieties/retained-preamble.tex lines 1-77), a verbatim excerpt of the native preamble.\par\medskip
\noindent\textit{Editorial note.} The label lemma-dimension-fibre-in-codim-1 is defined only in varieties.tex, so no cross-chapter restatement had to be weighed for this card.\par\medskip
\noindent\textit{Allowed standard prerequisites (editorial).}\begin{itemize}
\item Spec of the local ring at y is the set of points specializing to y. Schemes, Lemma schemes-lemma-specialize-points, invoked at proof lines 3606-3608. Supplies the reduction of the base to the local ring at y.
\item Dimension formula for a morphism of schemes. Morphisms, Lemma morphisms-lemma-dimension-formula, invoked at proof lines 3626-3627. Supplies the inequality from which the upper bound on the fibre dimension is read off.
\item Dimension of a fibre at a point as a sum of the local ring dimension and a transcendence degree. Morphisms, Lemma morphisms-lemma-dimension-fibre-at-a-point, invoked at proof lines 3635-3636 and 3642-3646. Supplies the identification of the fibre dimension with the two expressions compared by the dimension formula.
\item Dimension of fibres is bounded below by the dimension at the generic point of the component. Morphisms, Lemma morphisms-lemma-openness-bounded-dimension-fibres, invoked at proof lines 3647-3649. Supplies the reverse inequality dim\_x(X\_y) >= d.
\end{itemize}\medskip
\ifdefined\hypertarget\hypertarget{sourcepointer-1}{}\fi
\noindent\textit{Source reference.} Original lemma 02FX: sources/stacks-chow/tags.txt, source lines 3208--3208. Pointer outside the counted span, invoked once for context; the load-bearing argument of the counted statement is the author's proof reproduced in full in the proof span.\par\medskip
\ifdefined\hypertarget\hypertarget{sourcepointer-2}{}\fi
\noindent\textit{Source reference.} Original lemma 02JU: sources/stacks-chow/tags.txt, source lines 3349--3349. Pointer outside the counted span, invoked once for context; the load-bearing argument of the counted statement is the author's proof reproduced in full in the proof span.\par\medskip
\ifdefined\hypertarget\hypertarget{sourcepointer-3}{}\fi
\noindent\textit{Source reference.} Original lemma 02FZ: sources/stacks-chow/tags.txt, source lines 3210--3210. Pointer outside the counted span, invoked once for context; the load-bearing argument of the counted statement is the author's proof reproduced in full in the proof span.\par\medskip
\ifdefined\hypertarget\hypertarget{sourcepointer-4}{}\fi
\noindent\textit{Source reference.} Original lemma 01J7: sources/stacks-chow/tags.txt, source lines 2004--2004. Pointer outside the counted span, invoked once for context; the load-bearing argument of the counted statement is the author's proof reproduced in full in the proof span.\par\medskip
\begin{lemma}
\label{lemma-dimension-fibre-in-codim-1}
Let $f : X \to Y$ be locally of finite type. Let $x \in X$ be a point
with image $y \in Y$ such that $\mathcal{O}_{Y, y}$ is Noetherian of
dimension $\leq 1$. Let $d \geq 0$ be an integer such that for every
generic point $\eta$ of an irreducible component of $X$ which contains
$x$, we have $\dim_\eta(X_{f(\eta)}) = d$. Then $\dim_x(X_y) = d$.
\end{lemma}

\begin{proof}
Recall that $\Spec(\mathcal{O}_{Y, y})$
is the set of points of $Y$ specializing to $y$, see
Schemes, Lemma \ref{schemes-lemma-specialize-points}.
Thus we may replace $Y$ by $\Spec(\mathcal{O}_{Y, y})$
and assume $Y = \Spec(B)$ where $B$ is a Noetherian
local ring of dimension $\leq 1$ and $y$ is the closed point.
We may also replace $X$ by an affine neighbourhood of $x$.

\medskip\noindent
Let $X = \bigcup X_i$ be the irreducible components of $X$ viewed
as reduced closed subschemes. If we can show each fibre $X_{i, y}$
has dimension $d$, then $X_y = \bigcup X_{i, y}$ has dimension $d$ as
well. Thus we may assume $X$ is an integral scheme.

\medskip\noindent
If $X \to Y$ maps the generic point $\eta$ of $X$ to $y$, then $X$
is a scheme over $\kappa(y)$ and the result is true by assumption.
Assume that $X$ maps $\eta$ to a point $\xi \in Y$ corresponding
to a minimal prime $\mathfrak q$ of $B$ different from $\mathfrak m_B$.
We obtain a factorization $X \to \Spec(B/\mathfrak q) \to \Spec(B)$.
By the dimension formula
(Morphisms, Lemma \ref{morphisms-lemma-dimension-formula})
we have
$$
\dim(\mathcal{O}_{X, x}) + \text{trdeg}_{\kappa(y)} \kappa(x) \leq
\dim(B/\mathfrak q) + \text{trdeg}_{\kappa(\mathfrak q)}(R(X))
$$
We have $\dim(B/\mathfrak q) = 1$. We have
$\text{trdeg}_{\kappa(\mathfrak q)}(R(X)) = d$ by
our assumption that $\dim_\eta(X_\xi) = d$, see
Morphisms, Lemma \ref{morphisms-lemma-dimension-fibre-at-a-point}.
Since $\mathcal{O}_{X, x} \to \mathcal{O}_{X_s, x}$ has a kernel
(as $\eta \mapsto \xi \not = y$) and since $\mathcal{O}_{X, x}$
is a Noetherian domain we see that
$\dim(\mathcal{O}_{X, x}) > \dim(\mathcal{O}_{X_y, x})$.
We conclude that
$$
\dim_x(X_s) =
\dim(\mathcal{O}_{X_s, x}) + \text{trdeg}_{\kappa(y)} \kappa(x) \leq d
$$
(Morphisms, Lemma \ref{morphisms-lemma-dimension-fibre-at-a-point}).
On the other hand, we have
$\dim_x(X_s) \geq \dim_\eta(X_{f(\eta)}) = d$ by
Morphisms, Lemma \ref{morphisms-lemma-openness-bounded-dimension-fibres}.
\end{proof}

\end{document}
